\documentclass{ou}
\begin{document}
\thispagestyle{empty}
\begin{center}
\shortstack{\FGAVO}\vspace{4\baselineskip}
\textbf{РЕФЕРАТ}\\
\setstretch{2}
Информационные системы и~технологии\\
\setstretch{1}
Информационные технологии и информационные системы\vspace{2\baselineskip}\end{center}
\begin{flushleft}
\doublespacing
\begin{tabularx}{\textwidth}{@{} p{0.3\textwidth} l l @{}}
Руководитель & ст. преподаватель & И. Р. Нигматуллин\\
Студенты     &                   &                  \\
\multirow{-2.27}{\hsize}{\singlespacing Понятие и состав информационных технологий как элемента цифровой среды}
             & ГФ25-01Б, 152537671 & А. Е. Корикова\\
             & ГФ25-01Б, 152536528 & Д. А. Васильева\\
             & ГФ25-01Б, 152540016 & И. Н. Морозова\\
\multirow{-2.27}{\hsize}{\singlespacing Структура и функции информационных систем в современном обществе}
             & ГФ25-01Б, 152537684 & О. И. Бочанова\\
             & ГФ25-01Б, 152539986 & М. К. Гинтер\\
             & ГФ25-01Б, 152541471 & А. О. Безруких\\
             & ГФ25-01Б, 152536796 & А. В. Корюкова\\
\multirow{-2.27}{\hsize}{\singlespacing Взаимосвязь между информационными технологиями и~информационными системами}
             & ГФ25-01Б, 152540046 & А. Д. Балцевич\\
             & ГФ25-01Б, 152537659 & В. Е. Зайцева\\
             & ГФ25-01Б, 152536592 & А. Ю. Бирюкова\\
\end{tabularx}
\end{flushleft}
\vspace*{\fill}
\begin{center}
\city
\end{center}
\newpage
\thispagestyle{empty}
\setstretch{1}
Продолжение титульного~листа реферата по~теме «Информационные технологии и информационные системы»
\begin{flushleft}
\doublespacing
\begin{tabularx}{\textwidth}{@{} p{0.3\textwidth} l l @{}}
Студенты     &                     &         \\
\multirow{-2.27}{\hsize}{\singlespacing Классификация информационных систем по сфере применения и уровню автоматизации}
             & ГФ25-01Б, 152540753 & М. И. Бобылев\\
             & ГФ25-01Б, 152540765 & А. В. Гайворонская\\
             & ГФ25-01Б, 152539850 & Е. А. Аксаментова\\
\multirow{-2.27}{\hsize}{\singlespacing Роль информационных систем в управлении данными и поддержке принятия решений}
             & ГФ25-01Б, 152536115 & С. В. Идрисова\\
             & ГФ25-01Б, 152536056 & В. А. Карасова\\
             & ГФ25-01Б, 152539526 & К. Н. Аркадьева\\
             & ГФ25-01Б, 152540131 & Е. Кичко\\
\end{tabularx}
\end{flushleft}
\newpage
\tableofcontents
\newpage
\onehalfspacing
\newpage
\section{Понятие и состав информационных технологий как элемента цифровой среды}
Характерной чертой нашего времени являются развивающиеся процессы информатизации практически во всех сферах человеческой деятельности. Они привели к формированию новой информационной инфраструктуры, которая связана с~новым типом общественных отношений, с~новой реальностью, с~новыми информационными технологиями различных видов деятельности. 

Современному человеку необходимо знать информационные технологии, уметь успешно применять данные знания при решении как личностных, так и производственных задач повседневной жизни. Информационные технологии меняют нашу жизнь, культуру и профессиональную среду. Они формируют новую информационную среду, где знания становятся доступнее, коммуникации быстрее, а рабочие процессы -- эффективнее.
\subsection{Понятие информационных технологий}
Информационная технология -- совокупность методов, объединенная технологическим процессом, обеспечивающая сбор, хранение, обработку, вывод, 
распространение информации для~снижения трудоемкости процессов использования информационных ресурсов, повышения их надежности и~оперативности. Цель применения ИТ -- производство информации для~ее анализа человеком и принятия на его основе решения по выполнению какого-либо действия, а также снижение трудоемкости использования информационных ресурсов.

Методами информационных технологий являются методы сбора, передачи и обработки информации. Совокупность методов и производственных процессов определяет приемы, методы и мероприятия, регламентирующие проектирование и использование технических средств для обработки данных в предметной области.

Информационные технологии характеризуются следующими основными свойствами:
\begin{enumerate}[leftmargin=0em, itemindent=3.75em, nosep, label=\arabic*)]
\item Предметом (объектом) обработки (процесса) являются данные;
\item Целью процесса является получение информации;
\item Средствами осуществления процесса являются программные, аппаратные и программно-аппаратные вычислительные комплексы;
\item Процессы обработки данных разделяются на операции в соответствии с данной предметной областью (операция -- это совокупность элементарных действий, выполняемых на одном рабочем месте, которая приводит к~реализации определенной функции обработки данных);
\item Выбор управляющих воздействий на процессы осуществляется лицами, принимающими решение;
\item Критериями оптимизации процесса являются своевременность доставки информации пользователю, ее надежность, достоверность, полнота.
\end{enumerate}

Информационные технологии играют важную стратегическую роль, так как их применение позволило представить в формализованном виде, пригодном для практического использования, концентрированное выражение научных знаний и практического опыта для реализации и организации социальных процессов. Это привело к экономии затрат труда, времени, энергии, материальных ресурсов, необходимых для осуществления этих процессов.

Информационные технологии в настоящее время получили очень широкое применение. Они используются практически во всех сферах человеческой деятельности. В зависимости от рода предметной области выделяют три класса информационных технологий:
\begin{itemize}[label=-, leftmargin=0em, itemindent=3.35em, nosep]
\item Глобальные информационные технологии позволяют использовать информационные ресурсы общества в целом;
\item Базовые информационные технологии применяются для определенной области (производство, научные исследования, проектирование, обучение и~т.д.);
\item Конкретные информационные технологии реализуют обработку данных при решении конкретных функциональных задач пользователя (например, задачи планирования, учёта, анализа и т.д.).
\end{itemize}
\subsection{Понятие цифровой среды}
Цифровая среда представляет собой совокупность цифровых технологий, инфраструктуры и данных, которые обеспечивают взаимодействие и обмен информацией в цифровом формате. Это среда, в которой происходит создание, хранение, обработка и передача цифровой информации с использованием компьютерных и телекоммуникационных систем.

Цифровая среда играет ключевую роль в развитии современного общества и экономики. Она открывает новые возможности для бизнеса, образования, здравоохранения, государственного управления и многих других сфер.

Цифровизация позволяет повысить эффективность, скорость и качество предоставления услуг, а также расширить доступ к информации и знаниям. Кроме того, цифровая среда способствует развитию инноваций, создает новые рабочие места и стимулирует экономический рост.

В условиях глобализации развитие цифровой среды становится ключевым фактором конкурентоспособности и устойчивого развития как отдельных организаций, так и целых стран. Понимание роли и значения цифровой среды в современном мире является важным для успешной адаптации и эффективного использования цифровых технологий.
\subsection{Структура информационных технологий}
Структура информационных технологий представляет совокупность взаимосвязанных компонентов, совместная работа которых обеспечивает выполнение полного цикла обработки информации. Информационные технологии состоят из четырех основных компонентов: 
\begin{itemize}[label=-, leftmargin=0em, itemindent=3.35em, nosep]
\item Аппаратное обеспечение (hardware);
\item Программное обеспечение (software);
\item Информационные ресурсы;
\item Организационное обеспечение.
\end{itemize}

Аппаратное обеспечение включает в себя физическое оборудование, необходимое для функционирования информационных технологий.

Устройства обработки данных: центральные процессоры, графические процессоры, серверы, суперкомпьютеры.

Устройства хранения данных: жесткие диски, твердотельные накопители, системы хранения данных, ленточные библиотеки.

Устройства ввода/вывода информации: клавиатуры, манипуляторы, мониторы, принтеры, сканеры, сенсорные панели.

Сетевое и телекоммуникационное оборудование: маршрутизаторы, коммутаторы, модемы, точки доступа.

Компьютерные системы: персональные компьютеры, серверы.

Периферийные устройства: клавиатуры, мыши, мониторы, принтеры, сканеры, накопители данных.

Программное обеспечение -- это набор инструкций и приложений, которые управляют оборудованием и обеспечивают выполнение задач, определенных пользователем. Программное обеспечение включает в себя:

Системное программное обеспечение: операционные системы, драйверы устройств, утилиты, системы виртуализации. Обеспечивает базовое функционирование аппаратной части и управление ее ресурсами.

Прикладное программное обеспечение: программы, предназначенные для решения задач в конкретных предметных областях, такие как офисные пакеты, системы автоматизированного проектирования, бухгалтерские программы, корпоративные информационные системы.

Инструментальное программное обеспечение, или средства разработки: языки программирования, включая Python, Java, C++, интегрированные среды разработки, системы управления базами данных, средства тестирования. 

Информационные ресурсы составляют содержательную основу информационных технологий и включает все виды данных и знаний, циркулирующих в системе.

Структурированные данные: реляционные базы данных, электронные таблицы.

Неструктурированные и слабоструктурированные данные: текстовые документы, электронная почта, мультимедийный контент, изображения, аудио и видео.

Базы знаний и онтологии: формализованные модели, используемые в~экспертных системах и системах искусственного интеллекта.

Метаданные: данные, описывающие контекст, содержание и структуру других данных.

Сетевые и телекоммуникационные технологии: сохраняют объединение распределенных компонентов информационных технологий в единую систему.

Каналы связи: проводные, такие как оптоволоконные и медные, и~беспроводные, включая радиоканалы и спутниковые.

Сетевые протоколы: стандартизированные наборы правил для обмена данными, например, стек протоколов TCP/IP, HTTP/HTTPS, FTP, SMTP.

Сетевые технологии и архитектуры: технологии локальных, таких как Ethernet и Wi-Fi, и глобальных сетей, архитектуры «клиент-сервер», одноранговые сети.

Организационное обеспечение включает человеческие ресурсы и регламентирующую документацию, необходимые для эффективного использования и развития информационных технологий.

Персонал: ИТ-специалисты, такие как разработчики, администраторы, аналитики, руководители ИТ-направлений, конечные пользователи.

Организационно-распорядительная документация: регламенты, инструкции, политики информационной безопасности, стандарты оказания ИТ-услуг.

Методологии и системы управления: системы управления ИТ-проектами, управления конфигурациями, инцидентами и изменениями.

\subsection{Принципы реализации информационных технологий}
Основополагающими принципами реализации информационных технологий являются:
\begin{enumerate}[leftmargin=0em, itemindent=3.75em, nosep, label=\arabic*)]
\item Системный подход. Рассмотрение ИТ-систем как целостных структур, состоящих из взаимосвязанных элементов. Интеграция аппаратных, программных и информационных ресурсов в единую систему. Обеспечение согласованного взаимодействия компонентов системы;
\item  Модульность. Разделение ИТ-систем на независимые, заменяемые модули. Возможность расширения и модернизации системы путем добавления новых модулей;
\item Возможность наращивания мощности и функциональности ИТ-систем. Обеспечение гибкого распределения ресурсов для покрытия растущих потребностей;
\item Способность ИТ-систем быстро перестраиваться под изменяющиеся условия и требования. Использование адаптивных алгоритмов и механизмов обратной связи;
\item Надежность и отказоустойчивость. Обеспечение устойчивой работы ИТ-систем к сбоям, ошибкам и внешним воздействиям. Применение резервирования, репликации, балансировки нагрузки для повышения отказоустойчивости;
\item Безопасность. Защита ИТ-систем от несанкционированного доступа, утечки данных, вредоносных воздействий. Использование средств аутентификации, шифрования, межсетевых экранов, систем обнаружения вторжений.
\end{enumerate}
\subsection{Функционирование информационных технологий}
Функционирование информационных технологий представляет собой комплекс взаимосвязанных действий и процедур, обеспечивающих сбор, обработку, хранение, передачу и представление информации для решения различных задач в рамках информационных систем. Информационные технологии реализуют следующие основные функции:
\begin{enumerate}[leftmargin=0em, itemindent=3.75em, nosep, label=\arabic*)]
\item Сбор и ввод данных. Использование различных источников информации: первичные (наблюдения, эксперименты) и вторичные (документы, базы данных). Преобразование аналоговых данных в цифровую форму с~помощью сенсорных устройств;
\item Обработка и анализ данных. Применение математических методов обработки информации. Использование вычислительных технологий: облачные вычисления, высокопроизводительные системы, распределенные вычисления. Реализация аналитических и прогнозных моделей, визуализация данных;
\item Хранение и накопление данных. Использование различных типов носителей: магнитные, оптические, твердотельные, облачные. Организация баз данных и хранилищ данных для структурированного хранения информации. Обеспечение целостности, доступности и безопасности хранимых данных;
\item Передача и распространение данных Реализация каналов передачи данных: локальные, глобальные, мобильные сети;
\item Представление и визуализация данных. Использование графических, мультимедийных и других форматов для визуального представления информации. Создание интерфейсов и других средств эффективной визуализации данных;
\item Управление и администрирование. Организация информационных процессов и потоков в рамках ИТ-систем -- Применение методов и технологий управления жизненным циклом ИТ-систем. Контроль, оптимизация работы информационных технологий.
\end{enumerate}
\subsection{Связь информационных технологий и цифровой среды}
Цифровая среда представляет глобальное информационное пространство, формируемое совокупностью электронных средств и информационных, коммуникационных технологий, предназначенное для создания, хранения, обработки, передачи и использования информации в цифровой форме. В рамках этой среды информационные технологии выполняют роль системообразующего элемента, обеспечивающего ее функциональность, целостность.

Роль информационных технологий как элемента цифровой среды раскрывается через выполняемые ими ключевые функции:
\begin{itemize}[label=-, leftmargin=0em, itemindent=3.35em, nosep]
\item Формирование технологической инфраструктуры;
\item Обеспечение информационного взаимодействия;
\item Реализация цифровых сервисов и автоматизация процессов;
\item Управление информационными ресурсами.
\end{itemize}

Информационные технологии создают фундамент цифровой среды, который включает в себя:
\begin{enumerate}[leftmargin=0em, itemindent=3.75em, nosep, label=\arabic*)]
\item Аппаратно-техническую базу: серверы, центры обработки данных, системы хранения данных, клиентские устройства.
\item Программно-алгоритмическую базу: операционные системы, системы управления базами данных, платформенное и прикладное программное обеспечение.
\item Коммуникационную инфраструктуру: сети передачи данных, протоколы обмена, технологии беспроводной связи.
\end{enumerate}

Как отмечается в работах исследователей цифровой экономики, без этой инфраструктуры существование единой цифровой среды невозможно. Она является базисом для развертывания всех цифровых сервисов и платформ.

Информационные технологии реализуют процессы сбора, обработки, передачи и представления информации, обеспечивая коммуникацию между всеми субъектами цифровой среды: физическими лицами, организациями, государственными органами, машинами. Это взаимодействие осуществляется через стандартизированные протоколы и интерфейсы, что делает среду интегрированной и доступной.

На основе информационных технологий создаются и функционируют конкретные сервисы, составляющие пользовательский уровень цифровой среды. К ним относятся:
\begin{enumerate}[leftmargin=0em, itemindent=3.75em, nosep, label=\arabic*)]
\item Государственные информационные системы, например, портал «Госуслуги»;
\item Корпоративные информационные системы, такие как ERP и CRM;
\item Системы электронного документооборота;
\item Платформы для дистанционного образования и телемедицины.
\end{enumerate}

Согласно позиции в документах по цифровой трансформации, именно через эти сервисы цифровая среда оказывает непосредственное влияние на~социально-экономические процессы.

Информационные технологии предоставляют инструменты для работы с данными как стратегическим активом. Это включает технологии сбора, хранения, анализа, такие как Big Data и Data Mining, и защиты информации. Способность среды эффективно управлять большими объемами данных определяет ее ценность и потенциал для развития технологий искусственного интеллекта и аналитики.
\newpage
\section{Структура и функции информационных систем в~современном обществе}
Информационные системы (ИС) сегодня являются неотъемлемой частью современного общества, определяя ход его развития, функционирование экономики, управления, культуры и социальных связей. Их внедрение способствует автоматизации процессов, повышению эффективности, ускорению обмена информацией, а также созданию новых инновационных возможностей. Современная эпоха характеризуется развитием цифровых технологий, что делает информационные системы ключевыми инструментами для достижения целей развития и повышения уровня жизни населения. В~данном реферате представлено системное описание структуры и функций информационных систем, их эволюция, значение для общества и перспективы дальнейшего развития.
\subsection{Исторический аспект развития информационных систем}
Эволюция информационных систем прошла через несколько ключевых этапов, каждый из которых знаменовал переход к~более функциональным моделям:
\begin{itemize}[label=-, leftmargin=0em, itemindent=3.35em, nosep]
\item Этап автоматизации (1950–1960-е годы);
\item Этап развития встроенных и локальных систем (1970–1980-е годы);
\item Глобализация и развитие интернета (1990-е годы);
\item Постиндустриальный этап (2000–по настоящее время).
\end{itemize}
При этапе автоматизации появились первые электронные вычислительные машины (ЭВМ), таких как UNIVAC, IBM 701. Первичные системы предназначены для автоматизации бумажных процессов, бухгалтерский учет, расчет заработной платы, регистрация заказов.
Основа -- машинный язык программирования, увеличение скорости вычислений и первичные системы хранения данных.

Этап развития встроенных и локальных систем характеризуется внедрением локальных сетей, персональных компьютеров (PC). Создание систем управления предприятием (ERP), систем автоматизации производства, систем управления запасами и логистикой.
Формирование элементов корпоративных информационных систем, входящих в структуру бухгалтерии, маркетинга, планирования.

При глобализация и развитии интернета произошло создание Всемирной паутины (WWW) и массовое подключение к сети Интернет. Развитие электронной коммерции, электронных банков, систем обмена данными. Возникновение концепции «информационного общества».

При постиндустриальном этапе появились облачные технологии, мобильные устройства и развитие гигантских данных (Big Data). Широкое внедрение искусственного интеллекта (AI), системы машинного обучения и~автоматизированных решений. Внедрение технологий интернета вещей (IoT), робототехники, автоматизации умных городов.
\subsection{Современная структура информационных систем}
Современные ИС представляют собой сложные, многослойные, распределённые архитектуры, включающие в себя аппаратные, программные, информационные компоненты и коммуникационные сети:
\begin{enumerate}[leftmargin=0em, itemindent=3.75em, nosep, label=\arabic*)]
\item Аппаратное обеспечение (Hardware)
\item Программное обеспечение;
\item Базы данных и хранилища;
\item Коммуникационная инфраструктура;
\item Архитектурные модели.
\end{enumerate}

Аппаратное обеспечение включает в себя серверное оборудование и~дата-центры. Основные ядра хранения и обработки данных. Современные серверные фермы используют виртуализацию, энергоэффективные технологии. Клиентские устройства. Персональные компьютеры, ноутбуки, планшеты, смартфоны -- позволяют пользователю взаимодействовать с~системами. Мобильные устройства и IoT-устройства. Важнейшая часть современных ИС, обеспечивающая постоянную и удалённую доступность информации. Коммуникационные сети. Оптоволоконные линии, мобильные сети (4G, 5G), спутники и Wi-Fi -- обеспечивают мощную инфраструктуру передачи данных.

Программное обеспечение -- операционные системы. Windows, Linux, macOS -- управляют аппаратными ресурсами. Прикладные системы. Решения для автоматизации бизнес-процессов (ERP), управления взаимоотношениями с клиентами (CRM), систем документооборота, финансового учета. Облачные платформы. Amazon Web Services (AWS), Microsoft Azure, Google Cloud -- предоставляют услуги хранения, обработки и аналитики данных. Инструменты аналитики и искусственного интеллекта. Блоки машинного обучения, системы автоматического анализа и прогнозирования.

Реляционные базы данных (MySQL, Oracle, PostgreSQL) применяются для~структурированных, взаимосвязанных данных. NoSQL базы данных (MongoDB, Cassandra) -- для неструктурированных данных, обеспечивают масштабируемость. Большие хранилища данных и data lakes -- предназначены для литики, бизнес-отчетности, машинного обучения.

Коммуникационная инфраструктура:
\begin{enumerate}[leftmargin=0em, itemindent=3.75em, nosep, label=\arabic*)]
\item Локальные сети (LAN). Внутри предприятий и учреждений;
\item Глобальные сети (WAN). Взаимодействие филиалов и регионов;
\item Интернет и облачные платформы. Обеспечивают глобальный обмен информацией.
\end{enumerate}

Архитектурные модели:
\begin{itemize}[label=-, leftmargin=0em, itemindent=3.35em, nosep]
\item Клиент-серверные архитектуры;
\item Многоуровневые модели (трёхуровневая архитектура: презентационный слой, бизнес-логика, уровень данных);
\item Микросервисная архитектур;.
\item Облачные платформы и автоматизированные системы с использованием контейнеризации (Docker, Kubernetes).
\end{itemize}
\subsection{Основные функции современных информационных систем}
Информационные системы выполняют множество функций, которые подчинены целям повышения эффективности, автоматизации и управления:
\begin{enumerate}[leftmargin=0em, itemindent=3.75em, nosep, label=\arabic*)]
\item Обработка данных и информационный обмен. Автоматический сбор, структурирование, хранение и обработка данных. Обеспечение быстрого доступа к сведениям и обмена ими между участниками предприятия или общества. Использование систем интеграции данных (ETL-процессы) и~API-интерфейсов;
\item Автоматизация бизнес-процессов. Внедрение систем ERP, CRM, систем автоматического планирования, учета и контроля. Снижение затрат времени и человеческих ресурсов, повышение точности операций, снижение ошибок. Создание автоматизированных цепочек решений (workflow) для~цифровизации процессов;
\item Поддержка принятия управленческих решений (Decision Support Systems, DSS). Использование аналитической информации, прогнозов и~моделирования для руководства. Внедрение систем бизнес-аналитики (BI), дашбордов, отчетов. Обеспечение стратегического и оперативного управления;
\item Контроль и управление ресурсами. Мониторинг материальных активов, финансовых средств, человеческого капитала. Автоматизированное управление запасами, логистикой, финансами, кадрами. Использование систем мониторинга и сенсорных технологий.
\item Коммуникации и взаимодействие. Внутренние порталы, системы видеоконференций и мессенджеры. Электронная почта, системы совместной работы. Взаимодействие с внешними потребителями, поставщиками, государственными органами.
\item Обеспечение безопасности и защиты данных. Использование систем шифрования, аутентификация, контроль доступа. Кибербезопасность, системы предотвращения несанкционированных действий.
\end{enumerate}
\subsection{Значение информационных систем в современном обществе}
Информационные системы имеют огромное влияние на все стороны жизни современного человека и общества в целом:
\begin{enumerate}[leftmargin=0em, itemindent=3.75em, nosep, label=\arabic*)]
\item Влияние на экономику. Улучшение эффективности бизнеса и~конкуренции. Создание новых отраслей и рынков (финансовых технологий, электронной коммерции, умных устройств). Повышение прозрачности деловой среды, снижение издержек и автоматизация;
\item Развитие культуры, гуманитарных наук. Расширение образовательных онлайн-платформ, мультимедийных ресурсов. Обеспечение доступа к знаниям и культурным богатствам через цифровые библиотеки, музеи, наукоемкие проекты;
\item Обеспечение социальной интеграции и гражданского участия. Создание платформ для диалога между гражданами и органами власти. Инструменты для волонтерских, благотворительных и социальных инициатив;
\item Влияние на безопасность и управление гибкостью общества. Внедрение систем мониторинга и анализа угроз. Обеспечение мобильности, удаленной работы, электронных платежей, дистанционного образования;
\item Вызовы и риски. Угрозы кибербезопасности и защиты личных данных. Неравномерное распространение технологий, цифровой разрыв. Потеря рабочих мест из-за автоматизации и роботизации. Этические и~правовые аспекты использования ИС.
\end{enumerate}
\subsection{Современные тенденции и перспективы развития}
В будущем развитие информационных систем предполагает следующие ключевые направления:
\begin{enumerate}[leftmargin=0em, itemindent=3.75em, nosep, label=\arabic*)]
\item Интеграция искусственного интеллекта -- более интеллектуальные системы, способные самостоятельно обучаться и принимать решения. Облачные вычисления и инфраструктура как услуга (IaaS, PaaS, SaaS);
\item Интернет вещей (IoT) -- объединение устройств и объектов в единую сеть для автоматического сбора и обмена данными;
\item Блокчейн-технологии: обеспечение безопасности, децентрализации и~прозрачности транзакций;
\item Умные города и инфраструктуры: управление транспортом, энергопотреблением, безопасностью на базе ИС;
кибербезопасность и~защита данных -- постоянное совершенствование методов защиты информации в~свете увеличения масштабов киберугроз.
\end{enumerate}

Информационные системы в современном обществе представляют собой сложную и многослойную систему компонентов, объединённых в единую архитектуру, обеспечивающую автоматизацию, анализ, коммуникацию и~управление ресурсами. Их роль постоянно увеличивается, меняя структуру экономики, культуры и социальной жизни. В будущем развитие ИС будет связано с расширением возможностей искусственного интеллекта, роботизации и расширением цифровых технологий, что откроет новые горизонты для прогрессивного развития человеческого общества.
\newpage
\section{Взаимосвязь между информационными технологиями и информационными системами}
\subsection{Информационная технология}
Информационная технология -- совокупность четко определенных целенаправленных действий персонала по переработке информации на~компьютере.

Информационная технология является процессом, состоящим из~четко регламентированных правил выполнения операций, действий, этапов разной степени сложности над данными, хранящимися в компьютерах. Основная цель информационной технологии -- в результате целенаправленных действий по~переработке первичной информации получить необходимую для пользователя информацию.
\subsection{Информационная система}
Информационная система -- компьютерная система для~поддержки принятия решений и производства информационных продуктов, использующая компьютерную информационную технологию.

Информационная система представляет собой человеко-компьютерную систему обработки информации. Информационная система является средой, составляющими элементами которой является компьютеры, компьютерные сети, программные продукты, базы данных, люди, различного рода технические и программные средства, связи и т.д. Цель информационной системы -- организация хранения, обработки и передачи информации.
\subsection{Взаимосвязь}
Реализация функций информационной системы невозможна без знания ориентированной на нее информационной технологии. 

Информационная технология тесно связана с информационными системами, которые являются для нее основной средой. Информационные технологии охватывают ресурсы, необходимые для управления информацией, а информационные системы включают в себя данные, программы, аппаратное обеспечение и другие компоненты, которые позволяют автоматизировать процессы управления информацией. Более понятным языком, ИТ -- это инструментарий, а ИС -- это структура, которая этот инструментарий использует для достижения конкретных целей. ИТ и ИС не существуют друг без~друга в современном мире: ИТ предоставляют «двигатель» и~«инструменты», а~ИС -- это вся «машина» (включая водителя и пассажиров), которая использует эти инструменты для достижения определенных целей.

Взаимосвязь между ИТ и ИС:
\begin{enumerate}[leftmargin=0em, itemindent=3.75em, nosep, label=\arabic*)]
\item ИТ – основа ИС: любая информационная система строится на базе информационных технологий;
\item ИС задаёт требования к ИТ: информационная система определяет какие технологии нужны;
\item ИТ обеспечивают выполнение функций ИС;
\item ИТ без ИС -- просто набор технологий: просто наличие технологий ещё не создаёт систему. Компьютеры и программы не становятся информационной системой, пока они не включены в организованный процесс работы;
\item ИС развивается вслед за ИТ: когда появляются новые технологии (ИИ, Big Data, облака), ИС обновляются, чтобы использовать их. Например, раньше бухгалтерия велась в таблицах, но с появлением 1С стала полностью автоматизирована;
\item ИТ и ИС имеют общие цели: повышение эффективности, ускорение обработки данных и повышение уровня автоматизации.
\end{enumerate}

Таким образом, информационные технологии являются «строительными блоками» для информационных систем. В умелом сочетании информационных технологий -- управленческой и компьютерной -- залог успешной работы информационной системы.
\newpage
\section{Классификация информационных систем по сфере применения и уровню автоматизации}
\subsection{Понятие и сущность информационной системы}
Информационная система -- это совокупность взаимосвязанных методов, программно-технических средств, баз данных и персонала, предназначенная для сбора, передачи, обработки, хранения и выдачи информации. Основная задача информационных систем заключается в обеспечении пользователя актуальными и достоверными данными, необходимыми для принятия решений, анализа ситуации и управления различными процессами.

Структура любой информационной системы включает несколько ключевых компонентов. Во-первых, это аппаратное обеспечение, которое представляет собой вычислительную технику: персональные компьютеры, серверы, сетевое оборудование, устройства ввода и вывода информации. Во-вторых, программное обеспечение, обеспечивающее работу системы и её функциональные возможности: операционные системы, прикладные программы, сервисы и инструменты обработки данных.

Дополняют систему организационное обеспечение (инструкции, регламенты, правила работы), информационное обеспечение (базы данных, классификаторы, справочники, наборы знаний) и персонал, то есть специалисты, которые проектируют, обслуживают и используют систему.

Особенность любой информационной системы состоит в том, что она создаётся под конкретные задачи организации или сферы деятельности. Именно поэтому классификация по сфере применения позволяет выделить группы систем с учётом их функций.
\subsection{Классификация информационных систем по сфере применения}
По сфере применения информационные системы делятся на несколько категорий, каждая из которых ориентирована на решение специфических задач:
\begin{enumerate}[leftmargin=0em, itemindent=3.75em, nosep, label=\arabic*)]
\item Управленческие информационные системы (УИС);
\item Производственные информационные системы;
\item Финансово-экономические информационные системы;
\item Информационные системы в торговле;
\item Информационные системы в здравоохранении;
\item Информационные системы в государственном управлении;
\item Информационные системы в науке и исследовательской деятельности.
\end{enumerate}

Управленческие ИС используются для поддержки процессов управления на уровне организаций и предприятий. Они собирают и анализируют данные о~деятельности, формируют отчёты, помогают в планировании и контроле.

К числу распространённых систем относятся:
\begin{itemize}[label=-, leftmargin=0em, itemindent=3.35em, nosep]
\item MIS (Management Information Systems) -- системы информационной поддержки менеджмента, обеспечивающие доступ к данным для среднего управленческого звена;
\item EIS (Executive Information Systems) -- системы для руководителей высшего звена, предоставляющие стратегическую информацию;
\item ERP (Enterprise Resource Planning) -- системы планирования ресурсов предприятия, интегрирующие деятельность всех подразделений.
\end{itemize}

Использование подобных систем обеспечивает единое информационное пространство и повышает эффективность управленческих решений.

Производственные информационные системы нужны для~контроля процессов, обеспечения качества продукции и повышения эффективности производства.

К основным их типам относятся:
\begin{itemize}[label=-, leftmargin=0em, itemindent=3.35em, nosep]
\item SCADA-системы, обеспечивающие мониторинг и управление технологическими процессами;
\item MES-системы (Manufacturing Execution Systems), которые управляют производственными операциями в режиме реального времени;
\item APS-системы (Advanced Planning and Scheduling), выполняющие планирование и оптимизацию производственных задач.
\end{itemize}

Производственные ИС позволяют существенно снизить количество ошибок, оптимизировать ресурсы и повысить производственную дисциплину.

Финансово-экономические информационные системы применяются в~банках, страховых компаниях, бухгалтерии и инвестиционной сфере.

К ним относятся:
\begin{itemize}[label=-, leftmargin=0em, itemindent=3.35em, nosep]
\item Автоматизированные банковские системы;
\item Программы бухгалтерского учёта (например, 1С);
\item Аналитические инвестиционные платформы.
\end{itemize}

Основная задача данных систем -- точность финансовых операций, снижение рисков и автоматизация расчётных процессов.

В торговле информационные системы применяются для учёта товаров, анализа спроса, управления складом и автоматизации кассовых операций.

Среди наиболее распространённых:
\begin{itemize}[label=-, leftmargin=0em, itemindent=3.35em, nosep]
\item POS-системы, обеспечивающие работу касс;
\item CRM-системы, предназначенные для взаимодействия с клиентами;
\item Системы складского и логистического учёта.
\end{itemize}

Они обеспечивают оперативность обмена данными, улучшение обслуживания и прозрачность товарных потоков.

Медицинские ИС помогают вести электронные медицинские карты, хранить большие объёмы медицинских данных, организовывать работу врачей и поддерживать диагностические решения.

К ним относятся:
\begin{itemize}[label=-, leftmargin=0em, itemindent=3.35em, nosep]
\item Системы электронных медицинских карт;
\item Телемедицинские системы;
\item Системы поддержки диагностических решений (CDSS).
\end{itemize}

Использование таких систем улучшает и повышает качество диагностики и~эффективность функционирования медицинских учреждений.

Образовательные ИС применяются для автоматизации учебного процесса, управления учебным заведением и проведения дистанционного обучения.

В их числе:
\begin{itemize}[label=-, leftmargin=0em, itemindent=3.35em, nosep]
\item LMS (Learning Management Systems) -- системы дистанционного обучения;
\item Электронные журналы и дневники;
\item Системы тестирования и контроля знаний.
\end{itemize}

Такие решения делают образование более доступным и повышают эффективность педагогической деятельности.

Государственные ИС обеспечивают автоматизацию взаимодействия между государственными органами, гражданами и организациями.

Примеры:
\begin{itemize}[label=-, leftmargin=0em, itemindent=3.35em, nosep]
\item Порталы государственных услуг;
\item Системы электронного документооборота;
\item Государственные реестры и кадастры.
\end{itemize}

Они увеличивают прозрачность работы государственных органов и~сокращают время предоставления услуг.

Для научной деятельности создаются сложные информационные системы, позволяющие анализировать большие массивы данных, моделировать процессы и поддерживать публикационную активность.

К ним относятся:
\begin{itemize}[label=-, leftmargin=0em, itemindent=3.35em, nosep]
\item Системы научных публикаций;
\item Платформы обработки больших данных (Big Data);
\item Системы моделирования (например, MATLAB, ANSYS).
\end{itemize}
\subsection{Классификация информационных систем по уровню автоматизации}
По степени автоматизации информационные системы делятся на четыре основных вида:
\begin{enumerate}[leftmargin=0em, itemindent=3.75em, nosep, label=\arabic*)]
\item Ручные информационные системы;
\item Автоматизированные информационные системы;
\item Автоматические информационные системы;
\item Интеллектуальные информационные системы.
\end{enumerate}

Ручные системы основаны на обработке информации человеком без~использования компьютерных технологий. В настоящее время они встречаются редко, но могут применяться при отсутствии технических средств. Для них характерны низкая скорость обработки данных, высокая вероятность ошибки и недостаточная надёжность.

В автоматизированных системах значительная часть операций выполняется техническими средствами, однако человек продолжает принимать основные решения и контролировать процесс. Такие системы позволяют автоматизировать рутинные операции, но требуют подготовки персонала и~постоянного участия специалистов.

Автоматические ИС функционируют без участия человека: они самостоятельно собирают данные, обрабатывают их, принимают решения и формируют результаты. Участие человека требуется только в аварийных ситуациях или для общего контроля. К ним относятся системы автоматического регулирования, охранные системы, автономные комплексы мониторинга.

Интеллектуальные ИС основаны на методах искусственного интеллекта и анализе больших объёмов данных. Они способны самостоятельно выявлять закономерности, прогнозировать события, адаптироваться к изменениям и обучаться.
Примеры: рекомендательные сервисы, экспертные системы, интеллектуальные роботы и системы прогнозирования.

Анализ классификаций показывает, что выбор вида информационной системы зависит от характера задач и требований к автоматизации.

В здравоохранении использование полностью автоматических систем недопустимо при диагностике, поэтому преобладают автоматизированные и интеллектуальные решения, где человек остаётся ключевым звеном. В~производстве многие процессы могут выполнять автоматические системы, так как они требуют высокой точности и скорости реакции. В государственном управлении важно сохранять управленческий контроль, поэтому здесь широко применяются автоматизированные ИС.

В научной сфере, где важно анализировать большие массивы данных, активно используются интеллектуальные системы, способные выполнять сложную обработку информации. Таким образом, степень автоматизации напрямую связана с требованиями к надёжности, рискам и характером деятельности.
\newpage
\section{Роль информационных систем в управлении данными и поддержке принятия решений}
В современном мире, характеризующемся экспоненциальным ростом объемов данных, информационные системы (ИС) играют критически важную роль в эффективном управлении информацией и поддержке принятия обоснованных решений. Они позволяют организациям собирать, хранить, обрабатывать и анализировать данные, превращая их в ценные знания, необходимые для стратегического и оперативного управления.
\subsection{Информационные системы как инструмент управления данными}
Информационные системы представляют собой интегрированный набор компонентов, предназначенных для сбора, хранения, обработки и распространения информации. В контексте управления данными, ИС обеспечивают:
\begin{enumerate}[leftmargin=0em, itemindent=3.75em, nosep, label=\arabic*)]
\item Централизованное хранение данных: ИС позволяют консолидировать данные из различных источников в едином хранилище, обеспечивая их целостность и доступность. (Laudon \verb|&| Laudon, 2018);
\item Обеспечение качества данных: ИС включают механизмы контроля качества данных, такие как валидация, проверка на дублирование и~исправление ошибок, что гарантирует их достоверность и надежность. (O'Brien \verb|&| Marakas, 2011);
\item Защита данных: ИС обеспечивают защиту данных от потери и повреждения с помощью механизмов аутентификации, авторизации, шифрования и резервного копирования. (Whitman \verb|&| Mattord, 2017);
\item Управление метаданными: ИС позволяют управлять метаданными, т.е. данными о данных, что облегчает их поиск, понимание и использование. (Loshin, 2009)
\end{enumerate}

Примерами информационных систем, используемых для управления данными, являются системы управления базами данных (СУБД), хранилища данных (Data Warehouses) и озера данных (Data Lakes).
\subsection{Информационные системы как инструмент поддержки принятия решений}
Информационные системы предоставляют инструменты и методы для анализа данных и поддержки принятия решений на различных уровнях управления:
\begin{enumerate}[leftmargin=0em, itemindent=3.75em, nosep, label=\arabic*)]
\item Оперативный уровень: системы обработки транзакций (TPS) собирают и обрабатывают данные о текущих операциях, предоставляя информацию для контроля и управления повседневной деятельностью. (Turban et al., 2018);
\item Тактический уровень: системы управления информацией (MIS) предоставляют отчеты и аналитические данные, необходимые для планирования и контроля тактических задач. (Kenneth C. Laudon, Jane P. Laudon, 2018);
\item Стратегический уровень: системы поддержки принятия решений (DSS), системы бизнес-аналитики (BI) предоставляют инструменты для~анализа больших объемов данных, выявления тенденций и закономерностей, прогнозирования результатов и поддержки принятия стратегических решений. (Sharda et al., 2014).
\end{enumerate}
Информационные системы поддержки принятия решений используют различные методы анализа данных, такие как OLAP (Online Analytical Processing), Data Mining, Machine Learning и Predictive Analytics.

Примеры применения информационных систем в управлении данными и~поддержке принятия решений:
\begin{enumerate}[leftmargin=0em, itemindent=3.75em, nosep, label=\arabic*)]
\item Финансы: ИС используются для управления финансовыми данными, автоматизации бухгалтерского учета, анализа финансовых показателей и~прогнозирования финансовых результатов;
\item Маркетинг: ИС используются для управления данными о клиентах, анализа рынка, разработки маркетинговых стратегий и оценки эффективности маркетинговых кампаний;
\item Производство: ИС используются для управления производственными процессами, планирования производства, контроля качества и оптимизации логистики;
\item Здравоохранение: ИС используются для управления медицинскими данными, поддержки принятия клинических решений, мониторинга состояния пациентов и улучшения качества медицинского обслуживания;
\item Логистика: ИС используются для отслеживания грузов, оптимизации маршрутов, управления складскими запасами и повышения эффективности цепочки поставок.
\end{enumerate}

Информационные системы играют ключевую роль в управлении данными и поддержке принятия решений в современном бизнесе и других сферах деятельности. Они позволяют организациям эффективно собирать, хранить, обрабатывать и анализировать данные, превращая их в ценные знания, необходимые для повышения эффективности операционной деятельности, улучшения качества принимаемых решений и достижения стратегических целей. В условиях постоянно растущих объемов данных и усложняющихся задач управления, роль информационных систем будет только возрастать.
\newpage
{
\sloppy
\nocite{*}
\cleardoublepage
\addcontentsline{toc}{section}{Список использованных источников}
\bibliographystyle{plainurl}
\bibliography{report1}
}
\end{document}
